% Laget automatisk av elevark/lag_elevark.py fra opplegg/repetisjon-mekanikk-2.html. Endringer her blir overskrevet.
\documentclass[11pt]{article}
\usepackage{elevark}
\begin{document}
\elevarkhode{Fysikk 2 · Sirkelbevegelse}{Repetisjon: kast og sirkelbevegelse}{https://lektorbo.no/opplegg/repetisjon-mekanikk-2.html}{Opplegget på nett}

I Fysikk 2 blir bevegelsen todimensjonal. Nøkkelen er å dele vektorene i komponenter og behandle x- og y-retningen hver for seg. I sirkelbevegelse peker akselerasjonen inn mot sentrum. Her er en oversikt og blandede oppgaver med løsningsforslag.

\section*{Kast og sirkelbevegelse side om side}

\begin{center}\small
\begin{tabularx}{\linewidth}{YYY}
\toprule
\textbf{} & \textbf{Kast (uten luftmotstand)} & \textbf{Jevn sirkelbevegelse} \\
\midrule
Akselerasjon & $g$, rett ned, konstant & $v^{2}/r$, inn mot sentrum, retningen endres hele tiden \\
Kraftsum & tyngden mg & sentripetalkraften $mv^{2}/r$, summen av de virkelige kreftene inn mot sentrum \\
Fart & $v_{x}$ konstant, $v_{y}$ endres & konstant størrelse, retningen endres \\
Banen & parabel & sirkel \\
\bottomrule
\end{tabularx}
\end{center}

Sentripetalkraften er ikke en ny kraft. Det er navnet på kraftsummen inn mot sentrum. Den kan være snorkraft, friksjon, normalkraft, tyngde eller en sum av flere krefter.

\section*{Formler}

\begin{itemize}
\item Kast: $x = v_{0x}t$ og $y = y_{0} + v_{0y}t - \tfrac{1}{2}gt^{2}$, med $v_{0x} = v_{0} \cos  α$ og $v_{0y} = v_{0} \sin  α$. $v_{x} = v_{0x}$ og $v_{y} = v_{0y} - gt$.
\item Sirkelbevegelse: $v = 2πr/T = ωr$ og $a = v^{2}/r = ω^{2}r = 4π^{2}r/T^{2}$.
\item Dosert sving uten friksjon: $\tan  θ = v^{2}/(gr)$. Flat sving: $μmg \geq mv^{2}/r$.
\item Loddrett sirkel: på toppen $mg + N = mv^{2}/r$, i bunnen $N - mg = mv^{2}/r$. Minste fart på toppen: $v = \sqrt{gr}$.
\item Sirkelbane rundt en planet: $GMm/r^{2} = mv^{2}/r$ gir $v = \sqrt{GM/r}$.
\end{itemize}

Konstanter: $g = 9{,}81\,\mathrm{m/s^{2}}$, $G = 6{,}67 \cdot 10^{-11}\,\mathrm{N\,m^{2}/kg^{2}}$, GM for jorda = $3{,}986 \cdot 10^{14}\,\mathrm{m^{3}/s^{2}}$ og jordradius 6371~km.

\section*{Oppgaver}

Se bort fra luftmotstand hvis ikke noe annet står.

\begin{oppgaveliste}
\oppgave \fase{forklar}{Begrep}Er påstandene sanne eller usanne? Begrunn svaret.

a) På toppen av et skrått kast er akselerasjonen null.\newline b) En bil som kjører gjennom en sving med konstant fart, har akselerasjon.\newline c) Sentrifugalkraften presser deg utover i en karusell.\newline d) En kule som slippes, og en som skytes ut vannrett fra samme høyde, treffer bakken samtidig.\newline e) En satellitt i bane er vektløs fordi det ikke er tyngdekraft der.
\svar
\oppgave \fase{test}{Regn}En ball triller utfor en bordkant 0,90~m over gulvet med farten 2,5~m/s. Hvor lang tid bruker den, hvor langt fra bordet lander den, og med hvilken fart og vinkel treffer den gulvet? Se også Horisontalt kast med Eulers metode.\footnote{\url{lektorbo.no/programmering/python-horisontalt-kast.html}}
\svar
\oppgave \fase{test}{Regn}En ball sparkes fra bakken med 20~m/s, 35° over vannrett. Finn tiden i lufta, rekkevidden og den største høyden. Hvilken annen vinkel gir samme rekkevidde? Sjekk med simuleringen.\footnote{\url{lektorbo.no/sim/skratt-kast.html}}
\svar
\oppgave \fase{test}{Regn}Et lodd på 5,0~kg henger midt på en snor som er festet i to punkter i samme høyde. Hver halvdel av snora danner 30° med vannrett. Finn snorkraften. Hva skjer med snorkraften når snora strammes slik at vinkelen blir mindre? Sjekk med simuleringen.\footnote{\url{lektorbo.no/sim/lodd-i-snorer.html}}
\svar
\oppgave \fase{test}{Regn}En bil kjører i en flat sving med radius 50~m. Friksjonstallet er 0,70. Hva er den største farten bilen kan ha uten å skli? Hva blir den på glatt føre med $μ = 0{,}15$? Se Sirkelbevegelse.\footnote{\url{lektorbo.no/sim/sirkelbevegelse.html}}
\svar
\oppgave \fase{test}{Regn}En berg-og-dal-bane har en loddrett loop med radius 8,0~m. Hva er den minste farten vogna kan ha på toppen? I bunnen av en dal med samme radius har en passasjer på 60~kg farten 20~m/s. Finn normalkraften. Hvor mange $g$ kjenner passasjeren? Sjekk med simuleringen.\footnote{\url{lektorbo.no/sim/loddrett-sirkel.html}}
\svar
\oppgave \fase{test}{Regn}En sving på en motorvei har radius 200~m og skal tas i 30~m/s uten hjelp fra friksjon. Hvor mye må svingen doseres? Sjekk med simuleringen.\footnote{\url{lektorbo.no/sim/dosert-sving.html}}
\svar
\oppgave \fase{utfordring}{Sammensatt}Den internasjonale romstasjonen går i en tilnærmet sirkelbane 420~km over bakken. Finn farten og omløpstiden. Hvor stor er tyngdeakselerasjonen der, sammenlignet med på bakken? Hvorfor svever astronautene likevel? Se også Satellittbane med Eulers metode.\footnote{\url{lektorbo.no/programmering/python-satellittbane.html}}
\svar
\end{oppgaveliste}

\end{document}
